Since the ring under consideration is an integral domain, we deduce the following relation, valid by substitution in every @@M0@@-ring, and for all elements @@M1@@ such that @@M2@@ for @@M3@@, @@M4@@ for @@M5@@;
\sgaStableTarget{sga6.e05.eq.5-4-2}@@M6@@

\sgaStableTarget{sga6.e05.para.5-5}5.5. Now let @@M7@@ be any @@M8@@-ring, and let @@M9@@ be an element of @@M10@@ such that @@M11@@ for @@M12@@. We shall associate to @@M13@@ a new @@M14@@-ring @@M15@@.

We first define a new multiplication on @@M16@@ by putting, for @@M17@@:
@@M18@@

The ring @@M19@@ with this multiplication is no longer unital; we therefore consider the ring @@M20@@, obtained by adjoining a unit to @@M21@@. Its operations are thus defined by:
@@M22@@
We regard @@M23@@ as embedded in this ring by @@M24@@. We then define a pre-@@M25@@-ring structure by:
\sgaStableTarget{sga6.e05.eq.5-5-1}@@M26@@
This is a pre-@@M27@@-ring structure by (5.4.1). We shall denote by @@M28@@ the pre-@@M29@@-ring thus obtained. To see that it is a @@M30@@-ring, it suffices by substitution to give the proof in the case where @@M31@@ is the @@M32@@-ring generated over @@M33@@ by the generators @@M34@@, subject to the relations @@M35@@ for @@M36@@. We then define a homomorphism @@M37@@ by
